%% toc: P1 7 7.4,7.5,7.6,7.7
\section{Proofs for invariant arguments and changes of governing logic}
\label{sec:annex-p01-w02-p01-023-invariants-logic}

The notation is that of \cref{sec:p01-w02-p01-023-invariants-logic}. The
metalogic is classical.

\begin{annexitem}{annex:p01-w02-p01-023-invariant-principle}{thm:p01-w02-p01-023-invariant-principle}{The invariant principle}
\begin{proof}
\emph{Clause 1.} Let \(P\) be inductive and let \(s\) be reachable, the end
of a run \(s_0,\ldots,s_n\). We show \(s_k\in P\) for every \(k\leq n\) by
induction on \(k\). For \(k=0\), \(s_0\in I\subseteq P\) by (I1). If
\(k<n\) and \(s_k\in P\), then \(s_k\leadsto s_{k+1}\) by the definition of
a run, so \(s_{k+1}\in P\) by (I2). With \(k=n\) this gives \(s\in P\).
Hence \(\operatorname{Reach}\subseteq P\), so \(P\) is an invariant.

\emph{Clause 2.} (I1) for \(\operatorname{Reach}\): if \(s\in I\), the
one-term sequence \(s\) is a run of length \(0\) ending at \(s\), so
\(s\in\operatorname{Reach}\). (I2) for \(\operatorname{Reach}\): let
\(s\in\operatorname{Reach}\) and \(s\leadsto t\). Choose a run
\(s_0,\ldots,s_n\) with \(s_n=s\). Then \(s_0,\ldots,s_n,t\) satisfies
\(s_0\in I\) and each consecutive pair is a step, so it is a run of length
\(n+1\) ending at \(t\), and \(t\in\operatorname{Reach}\). Thus
\(\operatorname{Reach}\) is inductive. If \(P\) is any inductive property,
clause~1 gives \(\operatorname{Reach}\subseteq P\).

\emph{Clause 3.} If \(Q\) is an invariant, then
\(P:=\operatorname{Reach}\) is inductive by clause~2 and
\(P\subseteq Q\) by the definition of an invariant. Conversely, if \(P\) is
inductive and \(P\subseteq Q\), clause~1 gives
\(\operatorname{Reach}\subseteq P\subseteq Q\), so \(Q\) is an invariant.
\end{proof}
\end{annexitem}

\begin{annexitem}{annex:p01-w02-p01-023-variant}{thm:p01-w02-p01-023-variant}{Termination by a variant}
\begin{proof}
\emph{Clause 1.} Let \(s_0,\ldots,s_n\) be a run starting at \(s_0\). For
each \(k\leq n\) the initial segment \(s_0,\ldots,s_k\) is a run, so
\(s_k\) is reachable, and \(s_k\in P\) by
\cref{thm:p01-w02-p01-023-invariant-principle}, clause~1. For \(k<n\),
\(s_k\in P\) and \(s_k\leadsto s_{k+1}\), so \(\mu(s_{k+1})<\mu(s_k)\);
as both values are natural numbers, \(\mu(s_{k+1})\leq\mu(s_k)-1\). By
induction on \(k\), \(\mu(s_k)\leq\mu(s_0)-k\) for every \(k\leq n\). For
\(k=n\), \(0\leq\mu(s_n)\leq\mu(s_0)-n\), so \(n\leq\mu(s_0)\).

\emph{Clause 2.} Let \(L\) be the set of lengths of runs starting at
\(s_0\). The one-term sequence \(s_0\) is such a run, so \(0\in L\), and
\(L\subseteq\{0,\ldots,\mu(s_0)\}\) by clause~1. Hence \(L\) has a greatest
element \(N\); choose a run \(s_0,\ldots,s_N\). If some \(t\) satisfied
\(s_N\leadsto t\), then \(s_0,\ldots,s_N,t\) would be a run of length
\(N+1\) starting at \(s_0\), contrary to the choice of \(N\). So \(s_N\) is
terminal.
\end{proof}
\end{annexitem}

\begin{annexitem}{annex:p01-w02-p01-023-exit}{thm:p01-w02-p01-023-exit}{Closing at the exit}
\begin{proof}
\emph{Clause 1.} Let \(s\) be reachable and terminal. By
\cref{thm:p01-w02-p01-023-invariant-principle}, clause~1, \(s\in P\); and
\(s\in T\). Hence \(s\in P\cap T\subseteq Q\).

\emph{Clause 2.} Let \(\mu\) be a variant for \(P\) and \(s_0\in I\). By
\cref{thm:p01-w02-p01-023-variant}, clause~2, some run starting at \(s_0\)
ends at a terminal state \(s\). That state is reachable, since a run ends at
it, so \(s\in Q\) by clause~1. Every terminal state at which a run starting
at \(s_0\) ends is reachable for the same reason, so it lies in \(Q\) by
clause~1.
\end{proof}
\end{annexitem}

\begin{annexitem}{annex:p01-w02-p01-023-derivation-invariant}{thm:p01-w02-p01-023-derivation-invariant}{Invariants of a rule system}
\begin{proof}
The \emph{height} of a finite tree is \(0\) for a one-node tree and otherwise
one more than the largest height of the subtrees at the children of the
root. We prove, by strong induction on \(h\in\mathbb N\), that the root
label of every derivation of height \(h\) lies in \(P\).

Let \(D\) be a derivation of height \(h\), with root labelled \(u\) and
children of the root, in order, the roots of subtrees \(D_1,\ldots,D_k\)
labelled \(u_1,\ldots,u_k\) (here \(k=0\) when \(h=0\)). Each \(D_i\) is a
finite labelled tree in which every node satisfies the derivation condition,
because its nodes and their children are nodes and children of \(D\); so
\(D_i\) is a derivation, and its height is less than \(h\). By the induction
hypothesis \(u_i\in P\) for every \(i\leq k\); when \(k=0\) there is nothing
to check. The derivation condition at the root of \(D\) gives
\((\langle u_1,\ldots,u_k\rangle,u)\in\mathcal R\), and the hypothesis on
\(P\) gives \(u\in P\).

Every derivable item is the root label of some derivation, so it lies in
\(P\). If an item outside \(P\) were derivable, it would lie in \(P\); so no
item outside \(P\) is derivable.
\end{proof}
\end{annexitem}

\begin{annexitem}{annex:p01-w02-p01-023-limit}{thm:p01-w02-p01-023-limit}{Invariants at a limit stage}
\begin{proof}
\emph{Clause 1.} For each \(n\), the sequence \(s_0,\ldots,s_n\) is a run,
so \(s_n\) is reachable, and \(s_n\in P\) by
\cref{thm:p01-w02-p01-023-invariant-principle}, clause~1.

\emph{Clause 2.} Since \(s_n\leadsto s_{n+1}\), the hypothesis on steps
gives \(s_n\subseteq s_{n+1}\) for every \(n\), so
\(s_0\subseteq s_1\subseteq\cdots\). By clause~1 every \(s_n\) lies in
\(P\). Closure of \(P\) under increasing unions gives
\(\bigcup_{n\in\mathbb N}s_n\in P\).

\emph{Clause 3.} Let \(S\) consist of the finite subsets of \(\mathbb N\)
and the set \(\mathbb N\). For finite \(s\), the set
\(\mathbb N\setminus s\) is nonempty because \(\mathbb N\) is infinite, so
\(m_s:=\min(\mathbb N\setminus s)\) exists and \(s\cup\{m_s\}\) is a finite
subset of \(\mathbb N\), hence a state. The steps are exactly the pairs
\(s\leadsto s\cup\{m_s\}\) with \(s\) finite, so \(s\subseteq t\) whenever
\(s\leadsto t\). Let \(P\) be the set of finite states. (I1): \(\varnothing\)
is finite. (I2): if \(s\in P\) and \(s\leadsto t\), then
\(t=s\cup\{m_s\}\), which is finite. So \(P\) is inductive.

Let \(s_0,s_1,\ldots\) be an infinite sequence with \(s_0\in I\) and
\(s_n\leadsto s_{n+1}\) for every \(n\); the sequence
\(\{0,\ldots,n-1\}\), \(n\in\mathbb N\), shown below to satisfy these
conditions, is one, so such sequences exist. We show
\(s_n=\{0,\ldots,n-1\}\) by induction on \(n\). For \(n=0\),
\(s_0\in I=\{\varnothing\}\). If \(s_n=\{0,\ldots,n-1\}\), then \(s_n\) is
finite, the only step from \(s_n\) leads to \(s_n\cup\{m_{s_n}\}\), and
\(m_{s_n}=n\); so \(s_{n+1}=\{0,\ldots,n\}\). (The same computation shows
that \(\{0,\ldots,n-1\}\leadsto\{0,\ldots,n\}\) for every \(n\).) Every
\(m\in\mathbb N\) lies in \(s_{m+1}\), so
\(\bigcup_{n\in\mathbb N}s_n=\mathbb N\), which is a state and is not
finite. Thus every stage lies in \(P\) but the union does not, and \(P\) is
not closed under increasing unions.
\end{proof}
\end{annexitem}

\begin{annexitem}{annex:p01-w02-p01-023-horn-chaining}{thm:p01-w02-p01-023-horn-chaining}{Correctness of forward chaining}
\begin{proof}
Every \(h\) occurring in a clause \((B,h)\in\Phi\) is an atom occurring in
\(\Phi\) or is \(\bottom\), so \(h\in A\cup\{\bottom\}\), and
\(M\cup\{h\}\) is a state whenever \(M\) is. Since \(A\) is a set of atoms,
\(\bottom\notin A\) and \(\lvert A\cup\{\bottom\}\rvert=\lvert A\rvert+1\).

\emph{Clause 1.} (I1): no member of \(\varnothing\) needs to be true, so
\(\varnothing\in P_\Phi\). (I2): let \(M\in P_\Phi\) and
\(M\leadsto M\cup\{h\}\) by a clause \((B,h)\in\Phi\) with \(B\subseteq M\)
and \(h\notin M\). Let \(v\) be a valuation satisfying every clause of
\(\Phi\). Since \(M\in P_\Phi\), every member of \(M\) is true under \(v\);
since \(B\subseteq M\), every member of \(B\) is true under \(v\). As \(v\)
satisfies \((B,h)\) and makes no member of \(B\) false, \(h\) is an atom
that \(v\) makes true. Hence every member of \(M\cup\{h\}\) is true under
\(v\). As \(v\) was arbitrary, \(M\cup\{h\}\in P_\Phi\).

\emph{Clause 2.} Every state \(M\) is a subset of the finite set
\(A\cup\{\bottom\}\), so \(\mu_\Phi(M)\in\mathbb N\). If
\(M\leadsto M\cup\{h\}\), then \(h\notin M\), so
\(\lvert M\cup\{h\}\rvert=\lvert M\rvert+1\) and
\(\mu_\Phi(M\cup\{h\})=\mu_\Phi(M)-1<\mu_\Phi(M)\). So \(\mu_\Phi\) is a
variant for \(P_\Phi\). By \cref{thm:p01-w02-p01-023-variant}, every run
starts at \(\varnothing\) and has length at most
\(\mu_\Phi(\varnothing)=\lvert A\rvert+1\), and some run ends at a terminal
state.

\emph{Clause 3.} \(M\) is terminal if and only if no clause
\((B,h)\in\Phi\) has \(B\subseteq M\) and \(h\notin M\), that is, if and
only if \(h\in M\) for every \((B,h)\in\Phi\) with \(B\subseteq M\), which
is closure.

\emph{Clause 4.} Let \(M\) be reachable and terminal. By clause~1 and
\cref{thm:p01-w02-p01-023-invariant-principle}, clause~1,
\(M\in P_\Phi\).

Suppose \(\bottom\in M\). If a valuation \(v\) satisfied every clause of
\(\Phi\), then, as \(M\in P_\Phi\), every member of \(M\) would be true under
\(v\), including \(\bottom\); but \(\bottom\) is true under no valuation. So
no valuation satisfies every clause of \(\Phi\).

Suppose \(\bottom\notin M\). Let \(v_M\) make an atom true exactly when it
belongs to \(M\). Let \((B,h)\in\Phi\). If \(B\not\subseteq M\), some
\(b\in B\) is not in \(M\), so \(v_M\) makes \(b\) false and satisfies
\((B,h)\). If \(B\subseteq M\), then \(h\in M\) because \(M\) is closed
(clause~3); as \(\bottom\notin M\), \(h\) is an atom in \(M\), which
\(v_M\) makes true, so \(v_M\) satisfies \((B,h)\). Hence \(v_M\) satisfies
every clause of \(\Phi\).

The two cases give: if \(\bottom\in M\), no valuation satisfies \(\Phi\);
if \(\bottom\notin M\), some valuation does. Hence \(\bottom\in M\) if and
only if no valuation satisfies every clause of \(\Phi\), and the second case
supplies the stated satisfying valuation.
\end{proof}
\end{annexitem}

\begin{annexitem}{annex:p01-w02-p01-023-non-inductive}{thm:p01-w02-p01-023-non-inductive}{An invariant that must be strengthened}
\begin{proof}
\emph{Clause 1.} Let \(s_0,\ldots,s_n\) be a run. Then \(s_0=0\) and
\(s_{k+1}=s_k+2\) for \(k<n\), so \(s_k=2k\) by induction on \(k\), and the
run ends at \(2n\) with \(n\in\mathbb N\). Conversely, for \(m\in\mathbb N\)
the sequence \(0,2,\ldots,2m\) starts in \(I\), each term is the previous
one plus \(2\), and it ends at \(2m\). So
\(\operatorname{Reach}=\{2m:m\in\mathbb N\}\).

\emph{Clause 2.} If \(m=0\), then \(2m=0\neq1\); if \(m\geq1\), then
\(2m\geq2>1\). So \(\operatorname{Reach}\subseteq Q\).

\emph{Clause 3.} \(-1\in Q\) and \(-1\leadsto1\), because \(1=-1+2\); but
\(1\notin Q\). So (I2) fails for \(Q\).

\emph{Clause 4.} \(E=\{2m:m\in\mathbb Z\}\). (I1): \(0=2\cdot0\in E\).
(I2): if \(s=2m\) and \(s\leadsto t\), then \(t=2m+2=2(m+1)\in E\). If
\(1\) were in \(E\), then \(1=2m\) for an integer \(m\); but \(2m\leq0\)
for \(m\leq0\) and \(2m\geq2\) for \(m\geq1\). So \(1\notin E\), and
\(E\subseteq Q\).
\end{proof}
\end{annexitem}

\begin{annexitem}{annex:p01-w02-p01-023-classical-principles}{thm:p01-w02-p01-023-classical-principles}{Classical principles that are not intuitionistic}
\begin{proof}
\emph{Two facts used throughout.} (N) In every model and at every world
\(w\), \(w\forces\neg\varphi\) if and only if no \(v\geq w\) forces
\(\varphi\): by the implication clause, \(w\forces\varphi\to\bottom\) says
that every \(v\geq w\) forcing \(\varphi\) forces \(\bottom\), and no world
forces \(\bottom\). (C) For every valuation \(v\),
\(\widehat v(\varphi\to\bottom)\) is true exactly when
\(\widehat v(\varphi)\) is false, because \(\widehat v(\bottom)\) is false
and an implication is false exactly when its antecedent is true and its
consequent false (\cref{sec:prop-syntax-semantics-02}).

\emph{The models.} \(K_2\) has worlds \(w_0,w_1\), the order
\(w_0\leq w_0\), \(w_1\leq w_1\), \(w_0\leq w_1\), and the valuation with
\(V(p)=\{w_1\}\) and \(V(r)=\varnothing\) for every other atom \(r\), in
particular \(V(q)=\varnothing\). \(K_3\) has worlds \(w_0,w_1,w_2\), the
reflexive order with \(w_0\leq w_1\) and \(w_0\leq w_2\) added, and the
valuation with \(V(p)=\{w_1\}\), \(V(q)=\{w_2\}\) and
\(V(r)=\varnothing\) for every other atom \(r\). Each order is reflexive,
antisymmetric and transitive (the only pairs of distinct comparable worlds
have the form \(w_0\leq w_i\), and no chain of two such pairs exists), and
the world sets are nonempty. Each set \(V(r)\) is upward closed, because
\(w_1\) and \(w_2\) have no strictly greater worlds. So \(K_2\) and \(K_3\)
are intuitionistic Kripke structures in the sense of
\cref{sec:intuitionistic-logic-01} and models in the sense of
\cref{def:p08-ipl-completeness-semantics}.

\emph{Facts in \(K_2\).} (F1) \(w_1\forces p\) and \(w_0\nforces p\), by
the atomic clause. (F2) No world forces \(\neg p\): every world \(w\)
satisfies \(w\leq w_1\) and \(w_1\forces p\), so (N) fails. (F3) Every
world forces \(\neg\neg p\): by (F2) no \(v\geq w\) forces \(\neg p\), and
(N) applies. (F4) No world forces \(p\to q\): every world \(w\) has
\(w_1\geq w\) with \(w_1\forces p\) and \(w_1\nforces q\).

\emph{Non-derivability in clause 1.} By
\cref{thm:p08-ipl-completeness-soundness} with the empty premise set, if
\(\derives_{\mathsf{IPL}}\chi\) then every world of every model forces
\(\chi\). So it suffices to find a world that does not force \(\chi\).
\begin{enumerate}
  \item[(a)] In \(K_2\), \(w_0\nforces p\) by (F1) and \(w_0\nforces\neg p\)
        by (F2), so \(w_0\nforces p\lor\neg p\).
  \item[(b)] In \(K_2\), \(w_0\forces\neg\neg p\) by (F3) and
        \(w_0\nforces p\); as \(w_0\geq w_0\), the implication clause gives
        \(w_0\nforces\neg\neg p\to p\).
  \item[(c)] In \(K_2\), every \(v\geq w_0\) fails to force \(\neg p\) by
        (F2), so \(w_0\forces\neg p\to\neg\neg\neg p\) vacuously, while
        \(w_0\nforces\neg\neg p\to p\) by (b). As \(w_0\geq w_0\),
        \(w_0\nforces(\neg p\to\neg\neg\neg p)\to(\neg\neg p\to p)\).
  \item[(d)] In \(K_2\), no world forces \(p\to q\) by (F4), so every world
        forces \((p\to q)\to p\) vacuously; in particular \(w_0\) does, and
        \(w_0\nforces p\). As \(w_0\geq w_0\),
        \(w_0\nforces((p\to q)\to p)\to p\).
  \item[(e)] In \(K_3\), no world forces \(p\land q\): \(w_0\) forces
        neither atom, \(w_1\nforces q\) and \(w_2\nforces p\). By (N),
        \(w_0\forces\neg(p\land q)\). Also \(w_1\geq w_0\) forces \(p\), so
        \(w_0\nforces\neg p\) by (N); and \(w_2\geq w_0\) forces \(q\), so
        \(w_0\nforces\neg q\). Hence \(w_0\nforces\neg p\lor\neg q\), and as
        \(w_0\geq w_0\),
        \(w_0\nforces\neg(p\land q)\to(\neg p\lor\neg q)\).
\end{enumerate}

\emph{Classical validity in clause 1.} Let \(v\) be any valuation and use
(C) for every negation.
\begin{enumerate}
  \item[(a)] If \(\widehat v(p)\) is false, \(\widehat v(\neg p)\) is true;
        either way one disjunct of \(p\lor\neg p\) is true.
  \item[(b)] If \(\widehat v(\neg\neg p)\) is true, then
        \(\widehat v(\neg p)\) is false, so \(\widehat v(p)\) is true. So the
        antecedent is never true while the consequent is false, and
        \(\neg\neg p\to p\) is true.
  \item[(c)] Its consequent \(\neg\neg p\to p\) is true by (b), so the
        implication is true.
  \item[(d)] If \(\widehat v(p)\) is true, the formula is true. If
        \(\widehat v(p)\) is false, then \(\widehat v(p\to q)\) is true, so
        \(\widehat v((p\to q)\to p)\) is false, and the formula is true.
  \item[(e)] If \(\widehat v(\neg(p\land q))\) is true, then
        \(\widehat v(p\land q)\) is false, so \(\widehat v(p)\) or
        \(\widehat v(q)\) is false, so \(\widehat v(\neg p)\) or
        \(\widehat v(\neg q)\) is true, and the consequent is true.
\end{enumerate}
Since \(v\) was arbitrary, \(\entails_{\mathsf{CPL}}\chi\) for each of
(a)--(e).

\emph{Clause 2.} No world of any model forces both \(\neg\neg p\) and
\(\neg p\): such a world \(w\) would satisfy \(w\geq w\) and
\(w\forces\neg p\), so \(w\forces\bottom\) by the implication clause for
\(\neg\neg p=\neg p\to\bottom\). Hence
\(\{\neg\neg p,\neg p\}\entails_{\mathsf{IPL}}\bottom\) vacuously, and
\cref{thm:p08-ipl-strong-completeness-full} gives
\(\{\neg\neg p,\neg p\}\derives_{\mathsf{IPL}}\bottom\). By clause~1(b),
every valuation making \(\neg\neg p\) true makes \(p\) true, so
\(\{\neg\neg p\}\entails_{\mathsf{CPL}}p\). In \(K_2\), \(w_0\) forces
\(\neg\neg p\) (F3) and not \(p\) (F1), so
\(\{\neg\neg p\}\not\entails_{\mathsf{IPL}}p\), and
\cref{thm:p08-ipl-completeness-soundness} gives
\(\{\neg\neg p\}\notderives_{\mathsf{IPL}}p\). With
\(\Gamma=\{\neg\neg p\}\) and \(\varphi=p\), the premise judgement
\(\Gamma\cup\{\neg\varphi\}\derives_{\mathsf{IPL}}\bottom\) of the rule
holds and its conclusion \(\Gamma\derives_{\mathsf{IPL}}\varphi\) fails, so
the rule is not licensed in \(\mathsf{IPL}\).

\emph{Clause 3.} \(\entails_{\mathsf{CPL}}p\lor\neg p\) by clause~1(a). The
valuation making every atom false makes \(p\) false, so
\(\not\entails_{\mathsf{CPL}}p\). The valuation making every atom true makes
\(p\) true and hence \(\neg p\) false by (C), so
\(\not\entails_{\mathsf{CPL}}\neg p\).
\end{proof}
\end{annexitem}

\begin{annexitem}{annex:p01-w02-p01-023-repairs}{thm:p01-w02-p01-023-repairs}{Intuitionistic replacements and local repairs}
\begin{proof}
Each clause is proved by establishing the corresponding Kripke consequence
and then applying the completeness direction of
\cref{thm:p08-ipl-strong-completeness-full}; clauses (a) and (b) first use
its soundness direction (equivalently
\cref{thm:p08-ipl-completeness-soundness}) on their hypothesis. Persistence
of forcing, for single formulas and for premise sets, is
\cref{thm:p08-ipl-completeness-persistence}. Fact (N) of the preceding proof
is used for every negation: \(w\forces\neg\varphi\) if and only if no
\(v\geq w\) forces \(\varphi\). Every order is reflexive and transitive.
Throughout, \(w\) is an arbitrary world of an arbitrary model.

\emph{(a)} By soundness, \(\Gamma\cup\{\varphi\}\entails_{\mathsf{IPL}}\bottom\).
Let \(w\forces\Gamma\), and let \(v\geq w\) with \(v\forces\varphi\). By
persistence \(v\forces\Gamma\), so \(v\) forces every member of
\(\Gamma\cup\{\varphi\}\) and hence \(\bottom\), which no world forces. So
no \(v\geq w\) forces \(\varphi\), and \(w\forces\neg\varphi\) by (N). Thus
\(\Gamma\entails_{\mathsf{IPL}}\neg\varphi\), and completeness gives
\(\Gamma\derives_{\mathsf{IPL}}\neg\varphi\).

\emph{(b)} This is (a) with \(\neg\varphi\) in place of \(\varphi\), since
\(\neg\neg\varphi\) is \(\neg\varphi\to\bottom\).

\emph{(c)} Let \(v\geq w\) with \(v\forces\varphi\to\psi\); we show
\(v\forces\neg\psi\to\neg\varphi\). Let \(u\geq v\) with
\(u\forces\neg\psi\). If some \(x\geq u\) forced \(\varphi\), then
\(x\geq v\) by transitivity, so \(x\forces\psi\); and \(x\forces\neg\psi\)
by persistence; as \(x\geq x\), \(x\forces\bottom\), which is impossible.
So no \(x\geq u\) forces \(\varphi\), and \(u\forces\neg\varphi\) by (N).
Hence \(v\forces\neg\psi\to\neg\varphi\), and
\(w\forces(\varphi\to\psi)\to(\neg\psi\to\neg\varphi)\). This holds at
every world, so the formula is a Kripke consequence of the empty set, and
completeness gives
\(\derives_{\mathsf{IPL}}(\varphi\to\psi)\to(\neg\psi\to\neg\varphi)\).

\emph{(d), first formula.} Let \(w\forces\varphi\lor\neg\varphi\), so
\(w\forces\varphi\) or \(w\forces\neg\varphi\). Let \(v\geq w\) with
\(v\forces\neg\neg\varphi\); we show \(v\forces\varphi\). If
\(w\forces\varphi\), then \(v\forces\varphi\) by persistence. If
\(w\forces\neg\varphi\), then \(v\forces\neg\varphi\) by persistence, and
\(v\forces\neg\neg\varphi\) with \(v\geq v\) gives \(v\forces\bottom\),
which is impossible; so in this case no such \(v\) exists. Either way
\(w\forces\neg\neg\varphi\to\varphi\). So
\(\{\varphi\lor\neg\varphi\}\entails_{\mathsf{IPL}}\neg\neg\varphi\to\varphi\),
and completeness gives the derivability.

\emph{(d), second formula.} Let \(w\forces\psi\lor\neg\psi\). Let
\(v\geq w\) with \(v\forces\neg\psi\to\neg\varphi\), and let \(u\geq v\)
with \(u\forces\varphi\); we show \(u\forces\psi\). If \(w\forces\psi\),
then \(u\forces\psi\) by persistence, since \(u\geq w\). If
\(w\forces\neg\psi\), then \(v\forces\neg\psi\) by persistence; as
\(v\geq v\), \(v\forces\neg\varphi\); by persistence
\(u\forces\neg\varphi\); with \(u\forces\varphi\) and \(u\geq u\),
\(u\forces\bottom\), which is impossible; so in this case no such \(u\)
exists. Either way \(v\forces\varphi\to\psi\), and therefore
\(w\forces(\neg\psi\to\neg\varphi)\to(\varphi\to\psi)\). So
\(\{\psi\lor\neg\psi\}\entails_{\mathsf{IPL}}(\neg\psi\to\neg\varphi)\to(\varphi\to\psi)\),
and completeness gives the derivability.

\emph{(e)} Write \(X\) for \(\neg(p\land q)\to(\neg p\lor\neg q)\). Let
\(w\forces p\lor\neg p\), and let \(v\geq w\) with
\(v\forces\neg(p\land q)\); we show \(v\forces\neg p\lor\neg q\). If
\(w\forces\neg p\), then \(v\forces\neg p\) by persistence. If
\(w\forces p\), then \(v\forces p\) by persistence; if some \(u\geq v\)
forced \(q\), then \(u\forces p\) by persistence, so \(u\forces p\land q\),
while \(u\forces\neg(p\land q)\) by persistence, and \(u\geq u\) would give
\(u\forces\bottom\); so no \(u\geq v\) forces \(q\), and
\(v\forces\neg q\) by (N). In both cases \(v\forces\neg p\lor\neg q\). So
\(w\forces X\), that is,
\begin{equation*}
  \{p\lor\neg p\}\entails_{\mathsf{IPL}}X,
\end{equation*}
and completeness gives \(\{p\lor\neg p\}\derives_{\mathsf{IPL}}X\).

\emph{(f)} We show that every world forces \(\neg\neg X\); by (N) it
suffices that no \(v\geq w\) forces \(\neg X\). Suppose
\(v\geq w\) and \(v\forces\neg X\). Either some \(u\geq v\) forces \(p\),
or none does; this case distinction is made in the classical metalogic.
In the first case \(u\forces p\lor\neg p\), so \(u\forces X\) by the
displayed consequence in (e); also \(u\forces\neg X\) by persistence, and
\(u\geq u\) gives \(u\forces\bottom\), which is impossible. In the second
case \(v\forces\neg p\) by (N), so \(v\forces p\lor\neg p\) and
\(v\forces X\) by (e); with \(v\forces\neg X\) and \(v\geq v\),
\(v\forces\bottom\), which is impossible. So no such \(v\) exists, and
\(w\forces\neg\neg X\). Hence \(\neg\neg X\) is a Kripke consequence of the
empty set, and completeness gives \(\derives_{\mathsf{IPL}}\neg\neg X\).
\end{proof}
\end{annexitem}
